
%%%%%%%%%%%%%%%%%%%%%%%%%%%%%%%%%%%%%%%%
%           main body
%%%%%%%%%%%%%%%%%%%%%%%%%%%%%%%%%%%%%%%%


%-----------------------------------
%	TITLE PAGE
%-----------------------------------

\title[Convex Optimization]{\LARGE \bfseries 第四部分\quad 凸优化} % The short title appears at the bottom of every slide, the full title is only on the title page
\author[Exercises] % Your institution as it will appear on the bottom of every slide, may be shorthand to save space
{\Large \bfseries 习题}
% \author{Singular Value Decomposition} % Your name
% \date{\today} % Date, can be changed to a custom date
\date{}

%------------------------------------------------

\begin{document}


%------------------------------------------------

\begin{frame}

\titlepage % Print the title page as the first slide

\end{frame}

%------------------------------------------------

% \begin{frame}
% \frametitle{内容提要} % Table of contents slide, comment this block out to remove it
% \tableofcontents % Throughout your presentation, if you choose to use \section{} and \subsection{} commands, these will automatically be printed on this slide as an overview of your presentation
% \end{frame}

%-----------------------------------
%	PRESENTATION SLIDES
%-----------------------------------

%------------------------------------------------

\begin{frame}

{\color{red} \scriptsize 带$^*$号的题目可以不做}

\vspace{1em}

1. 验证以下集合是凸集
\begin{itemize}
\item[(1)] 平板，即形如$\{\mathbf{x} \in \mathbf{R}^n \ |\ a \le \mathbf{v}^T\mathbf{x} \le b \}$的集合
\item[(2)] 超矩形，即形如$\{\mathbf{x} \in \mathbf{R}^n \ |\ a_i \le x_i \le b_i, \; i=1,\ldots,n \}$的集合
\item[(3)] 楔形，即形如$\{\mathbf{x} \in \mathbf{R}^n \ |\ \mathbf{v}_1^T\mathbf{x} \le b_1, \; \mathbf{v}_2^T\mathbf{x} \le b_2 \}$的集合
\item[(4)] 距离一个集合比距离另一个集合近的点组成的集合，即
$$\{\mathbf{x} \in \mathbf{R}^n \ |\ \operatorname{dist}(\mathbf{x}, S_1) \le \operatorname{dist}(\mathbf{x}, S_2) \}$$
其中$S_1,S_2 \subset \mathbb{R}^n$, $\operatorname{dist}(\mathbf{x}, S) := \inf\limits_{\mathbf{y} \in S} \{ \Vert \mathbf{x} - \mathbf{y} \Vert_2 \}$
\item[(5)] 集合$\{\mathbf{x} \in \mathbf{R}^n \ |\ \mathbf{x} + S_2 \subset S_1 \}$, 其中$S_1,S_2 \subset \mathbb{R}^n$并且$S_1$是凸集，$\mathbf{x} + S_2 := \{ \mathbf{x} + \mathbf{y} \ |\ \mathbf{y} \in S_2 \}$
\end{itemize}

\end{frame}

%------------------------------------------------

\begin{frame}

2. 验证以下函数是凸函数
\begin{itemize}
\item[(1)] 指数和对数：$f(\mathbf{x}) = \ln(e^{x_1} + \cdots + e^{x_n})$
\item[(2)] 几何平均：$f(\mathbf{x}) = \left( \prod\limits_{i=1}^n x_i \right)^{1/n}$
\item[(3)] 对数行列式：$f(X) = \ln\det X$
\item[(4)] 凸函数的滑动平均（Running Average）
$$F(x) = \dfrac{1}{x}\int_0^{x} f(t) dt$$
其中$f$是一个定义域包含$\mathbb{R}_{\ge 0}$的凸函数
\end{itemize}

\vspace{1em}

3. 请证明，对于凸优化问题来说，局部最优解就是全局最优解。

\end{frame}

%------------------------------------------------

\begin{frame}

4. Logistic 模型的优化：设随机变量$X\in\{0,1\}$满足
$$P(X=1) = p(\mathbf{x}) = \dfrac{e^{\mathbf{a}^T\mathbf{x} + b}}{1 + e^{\mathbf{a}^T\mathbf{x} + b}}$$
其中$\mathbf{x} \in \mathbb{R}^n$是影响概率的向量变量，$\mathbf{a}, b$是已知参数。还假设向量变量$\mathbf{x}$受线性的不等式约束
$$\mathbf{f}_i^T\mathbf{x} \le g_i, \; i=1,\ldots,m$$
请将下列问题建模为凸优化问题
\begin{itemize}
\item[(1)] 选择$\mathbf{x}$以极大化$p(\mathbf{x})$
\item[(2)] 假设对所有可行的（即满足约束条件）的$\mathbf{x}$都有$\mathbf{c}^T\mathbf{x} + d > 0$. 选择$\mathbf{x}$以极大化$p(\mathbf{x})\cdot (\mathbf{c}^T\mathbf{x} + d)$
\end{itemize}

\end{frame}

%------------------------------------------------

\begin{frame}

5. 考虑标准凸优化问题
\begin{eqnarray*}
& \text{minimize} & f_0(\mathbf{x}) \\
& \text{subject to} & f_i(\mathbf{x}) \le 0, \quad i = 1,\ldots,m \\
& & h_j(\mathbf{x}) = 0, \quad j = 1,\ldots,p
\end{eqnarray*}
设$p^*$是最优值，$g(\boldsymbol\lambda, \boldsymbol\nu)$是对偶函数，请证明
$$p^* \ge g(\boldsymbol\lambda, \boldsymbol\nu)$$
对任意$\boldsymbol\lambda, \boldsymbol\nu$成立。

\end{frame}

%------------------------------------------------

\begin{frame}

6. 推导最小二乘法问题
\begin{eqnarray*}
& \text{minimize} & \Vert \mathbf{x} \Vert_2^2 (= \mathbf{x}^T\mathbf{x}) \\
& \text{subject to} & A\mathbf{x} = \mathbf{b}
\end{eqnarray*}
的对偶问题

\vspace{1em}

7. 解析中心（Analytic Centering）：考虑优化问题
$$\text{minimize} \quad -\sum\limits_{i=1}^m \ln(b_i - \mathbf{a}_i^T\mathbf{x})$$
推导其对偶问题

\end{frame}

%------------------------------------------------

\begin{frame}

8$^*$. 用KKT条件求解下列凸优化问题
\begin{eqnarray*}
& \text{minimize} & -\sum\limits_{i=1}^n \ln(\alpha_i + x_i) \\
& \text{subject to} & x_i \ge 0,\; i=1,\ldots,n \\
& & x_1+\cdots+x_n = 1
\end{eqnarray*}

\vspace{1em}

9. 请证明Newton减量$\lambda(\mathbf{x})$满足
\begin{align*}
\lambda(\mathbf{x}) & = \sup\limits_{\mathbf{v}^T H(f)(\mathbf{x}) \mathbf{v} = 1} \left( -\mathbf{v}^T \nabla f(\mathbf{x}) \right) \\
& = \sup\limits_{\mathbf{v} \neq \mathbf{0}} \dfrac{-\mathbf{v}^T \nabla f(\mathbf{x})}{(\mathbf{v}^T H(f)(\mathbf{x}) \mathbf{v})^{1/2}}
\end{align*}

\end{frame}

%------------------------------------------------

\begin{frame}

10. 求对数障碍函数$\phi(\mathbf{x}) := \sum\limits_{i=1}^m \left( -\ln(-f_i(\mathbf{x})) \right)$的梯度向量与黑塞矩阵

\vspace{1em}

11. 考虑对数障碍的近似问题
\begin{eqnarray*}
& \text{minimize} & tf_0(\mathbf{x}) + \phi(\mathbf{x}) \\
& \text{subject to} & A\mathbf{x} = \mathbf{b}
\end{eqnarray*}
令$\mathbf{x}^*(t)$为这个问题的最优解。求证
\begin{itemize}
\item[(1)] $\mathbf{x}^*(t)$是严格可行的，即满足
$$f_i(\mathbf{x}^*(t)) < 0, \; i = 1,\ldots,m; \quad A\mathbf{x}^*(t) = \mathbf{b}$$
\item[(2)] 存在$\widehat{\boldsymbol\nu} \in \mathbb{R}^p,$ 使得
$$\mathbf{0} = t \nabla f_0(\mathbf{x}^*(t)) + \nabla \phi(\mathbf{x}^*(t)) + A^T \widehat{\boldsymbol\nu}$$
\end{itemize}

\end{frame}

%------------------------------------------------

\begin{frame}

12$^*$. 继续考虑第11题中的问题。定义
$$\lambda_i^*(t) = - \dfrac{1}{tf_i(\mathbf{x}^*(t))}, \; i = 1,\ldots,m; \quad \boldsymbol\nu^*(t) = \widehat{\boldsymbol\nu} / t$$
请证明
\begin{itemize}
\item[(1)] $(\boldsymbol\lambda^*(t), \boldsymbol\nu^*(t))$是原问题的对偶可行解
\item[(2)] $f_0(\mathbf{x}^*(t)) - p^* \le \dfrac{m}{t}$
\end{itemize}

\end{frame}

%------------------------------------------------

% \begin{frame}

% \begin{block}{Vandermonde行列式的计算}
% \begin{enumerate}
% % \addtocounter{enumi}{0}
% \item 把第$n-1$行乘以$-a_1$加到第$n$行，再把第$n-2$行乘以$-a_1$加到第$n-1$行。按照这种顺序依次向上，直到用第$2$行乘以$-a_1$加到第$1$行。由行列式性质，$V_n$的值不变。
% \[
% V_n = \det \begin{pmatrix}
% 1 & 1 & \cdots & 1 \\
% 0 & a_2-a_1 & \cdots & a_n-a_1 \\
% \vdots & \vdots & & \vdots \\
% 0 & a_2^{n-2}(a_2-a_1) & \cdots & a_n^{n-2}(a_n-a_1) \\
% \end{pmatrix}
% \]
% \end{enumerate}
% \end{block}

% \end{frame}

%------------------------------------------------
% % closing page
% \begin{frame}
% \HUGE{\centerline{\bfseries {\color{gray} The} {\color{zkblue} End}}}

% \pause

% \vspace{1.2em}

% \Large{\centerline{\bfseries {\color{gray}Thank} \quad {\color{zkblue} You}}}

% \end{frame}

%----------------------------------------------------------------------------------------

\end{document}
